\documentclass[15pt]{extarticle}

\usepackage{amsmath}
\usepackage{amssymb}
\usepackage[
  paperwidth=13.333in,
  paperheight=7.5in,
  margin=0.55in
]{geometry}

\usepackage{xcolor}
\usepackage[most]{tcolorbox}
\usepackage{titlesec}
\usepackage{enumitem}
\usepackage{lmodern}

% ---------- Colors ----------
\definecolor{titleblue}{RGB}{20, 65, 120}
\definecolor{boxgray}{RGB}{248, 248, 248}
\definecolor{linegray}{RGB}{120, 120, 120}

% ---------- Section style ----------
\titleformat{\section}
  {\normalfont\Large\bfseries\color{titleblue}}
  {}
  {0pt}
  {}

% ---------- Spacing ----------
\setlength{\parindent}{0pt}
\setlength{\parskip}{0.20em}
\setlength{\abovedisplayskip}{0.35em}
\setlength{\belowdisplayskip}{0.35em}
\setlist[itemize]{leftmargin=1.2em, itemsep=0.12em, topsep=0.08em}

% ---------- Box style ----------
\tcbset{
  enhanced,
  colback=boxgray,
  colframe=linegray,
  boxrule=0.7pt,
  arc=2mm,
  left=8pt,
  right=8pt,
  top=6pt,
  bottom=6pt
}

\begin{document}
\large

% ============================================================
\newpage
\section*{1. Decision Problem}

\begin{tcolorbox}
\textbf{Goal:} Allocate the first limited vaccine supply to \textbf{one} of 10 cities to maximize expected future infection-burden reduction.
\end{tcolorbox}

\vspace{0.12cm}

\begin{minipage}[t]{0.47\textwidth}
\textbf{Epidemic state for city \(i\):}
\[
X_i(t)=
\begin{pmatrix}
S_i(t)\\
E_i(t)\\
I_i(t)\\
R_i(t)
\end{pmatrix},
\qquad i=1,\ldots,10.
\]

\textbf{Observed context:}
\[
C_i=\{I_i^0(t),R_i^0(t)\}_{t=1}^{30}.
\]

\textbf{Hidden future:}
\[
t=31,\ldots,60.
\]
\end{minipage}
\hfill
\begin{minipage}[t]{0.47\textwidth}
\textbf{Action space:}
\[
A=\{a_1,a_2,\ldots,a_{10}\},
\]
\[
a_i=\text{allocate vaccine to city }i.
\]

\textbf{Pre-intervention parameters:}
\[
\theta_i=(\beta_i,\sigma_i,\gamma_i).
\]

\[
\beta_i=\text{infection rate},\quad
\sigma_i=\text{progression rate},\quad
\gamma_i=\text{recovery rate}.
\]
\end{minipage}

\vspace{0.22cm}

\begin{tcolorbox}
This is a \textbf{decision-under-uncertainty} problem, not just a prediction problem:
\[
\text{choose }a_i
\quad\text{to maximize expected vaccine benefit.}
\]
\end{tcolorbox}


% ============================================================
\newpage
\section*{2. Pre-Intervention SEIR-SDE Model}

\textbf{Deterministic SEIR drift:}
\[
S \rightarrow E \rightarrow I \rightarrow R.
\]

\vspace{0.10cm}

\begin{minipage}[t]{0.47\textwidth}
\[
f_{S,i}= -\beta_i \frac{S_i I_i}{N_i},
\]
\[
f_{E,i}= \beta_i \frac{S_i I_i}{N_i}-\sigma_i E_i.
\]
\end{minipage}
\hfill
\begin{minipage}[t]{0.47\textwidth}
\[
f_{I,i}= \sigma_i E_i-\gamma_i I_i,
\]
\[
f_{R,i}= \gamma_i I_i.
\]
\end{minipage}

\vspace{0.18cm}

\textbf{Stochastic SEIR-SDE form:}
\[
\begin{pmatrix}
dS_i & dE_i & dI_i & dR_i
\end{pmatrix}
=
\begin{pmatrix}
f_{S,i} & f_{E,i} & f_{I,i} & f_{R,i}
\end{pmatrix}dt
+
dW_i(t)g_i(t,x).
\]

\[
dW_i(t)=
\begin{pmatrix}
dW_{1,i} & dW_{2,i} & dW_{3,i} & dW_{4,i}
\end{pmatrix}.
\]

\vspace{0.08cm}

\[
g_i(t,x)=
\begin{pmatrix}
-\sqrt{\beta_i \frac{S_i I_i}{N_i}}
& \sqrt{\beta_i \frac{S_i I_i}{N_i}}
& 0
& 0 \\[4pt]
0
& -\sqrt{\sigma_i E_i}
& \sqrt{\sigma_i E_i}
& 0 \\[4pt]
0
& 0
& -\sqrt{\gamma_i I_i}
& \sqrt{\gamma_i I_i} \\[4pt]
0
& 0
& 0
& 0
\end{pmatrix}
=
\begin{pmatrix}
g_{S,i} & g_{E,i} & g_{I,i} & g_{R,i}
\end{pmatrix}.
\]

\begin{tcolorbox}
The \(4\times4\) diffusion matrix is kept for consistency with the reference form. The fourth Wiener process does not create an additional epidemic transition because the fourth row is zero.
\end{tcolorbox}


% ============================================================
\newpage
\section*{3. Vaccine Counterfactual at Day 30}

\textbf{For candidate action \(a_i\), vaccine is applied to city \(i\) at day 30.}

\vspace{0.08cm}

\[
P_i=\min(VE\cdot Q,\;S_i^0(30)).
\]

\begin{tcolorbox}
Here, \(VE\cdot Q\) is the effective number of susceptible individuals protected by the available vaccine supply.
\end{tcolorbox}

\vspace{0.15cm}

\begin{minipage}[t]{0.47\textwidth}
\textbf{State intervention}

\[
S_i^{V}(30^+)=S_i^0(30)-P_i,
\]
\[
R_i^{V}(30^+)=R_i^0(30)+P_i.
\]

\[
E_i^{V}(30^+)=E_i^0(30),
\]
\[
I_i^{V}(30^+)=I_i^0(30).
\]
\end{minipage}
\hfill
\begin{minipage}[t]{0.47\textwidth}
\textbf{Parameter intervention}

\[
\beta_i^{V}=(1-\eta_i)\beta_i,
\qquad 0\leq \eta_i\leq 1.
\]

\[
\sigma_i^{V}=\sigma_i,
\qquad
\gamma_i^{V}=\gamma_i.
\]

\[
\theta_i^{V}=(\beta_i^{V},\sigma_i,\gamma_i).
\]
\end{minipage}

\vspace{0.25cm}

\begin{tcolorbox}
Vaccination reduces future infection through two mechanisms:
\[
S_i^{V}(30^+)<S_i^0(30),
\qquad
\beta_i^{V}<\beta_i.
\]
It reduces the susceptible population and lowers the effective transmission rate.
\end{tcolorbox}


% ============================================================
\newpage
\section*{4. Post-Vaccine Transmission and Dynamics}

\textbf{No-vaccine infection force:}
\[
\lambda_i^0(t)
=
\beta_i
\frac{S_i^0(t)I_i^0(t)}{N_i}.
\]

\textbf{Vaccine-counterfactual infection force:}
\[
\lambda_i^{V}(t)
=
\beta_i^{V}
\frac{S_i^{V}(t)I_i^{V}(t)}{N_i}.
\]

Since
\[
\beta_i^{V}=(1-\eta_i)\beta_i,
\]
we have
\[
\lambda_i^{V}(t)
=
(1-\eta_i)\beta_i
\frac{S_i^{V}(t)I_i^{V}(t)}{N_i}.
\]

\vspace{0.12cm}

\begin{tcolorbox}
Vaccination changes the effective transmission process:
\[
\lambda_i^{V}(t)
<
\lambda_i^0(t)
\]
when \(\eta_i>0\) and \(S_i^{V}(t)<S_i^0(t)\).
\end{tcolorbox}

\vspace{0.14cm}

\textbf{Post-vaccine drift terms:}
\[
f_{S,i}^{V}= -\beta_i^{V}\frac{S_i^{V}I_i^{V}}{N_i},
\qquad
f_{E,i}^{V}= \beta_i^{V}\frac{S_i^{V}I_i^{V}}{N_i}-\sigma_i E_i^{V},
\]
\[
f_{I,i}^{V}= \sigma_i E_i^{V}-\gamma_i I_i^{V},
\qquad
f_{R,i}^{V}= \gamma_i I_i^{V}.
\]

\begin{tcolorbox}
The intervention changes \(\beta_i\) into \(\beta_i^{V}\), but disease progression and recovery are unchanged:
\[
\sigma_i^{V}=\sigma_i,
\qquad
\gamma_i^{V}=\gamma_i.
\]
\end{tcolorbox}


% ============================================================
\newpage
\section*{5. Decision Rule and Evaluation}

\textbf{For each action \(a_i\), simulate two future trajectories from day 31 to day 60:}
\[
I_i^0(t)=\text{no-vaccine trajectory},
\qquad
I_i^V(t\mid a_i)=\text{vaccine-counterfactual trajectory}.
\]

\vspace{0.10cm}

\begin{minipage}[t]{0.47\textwidth}
\textbf{Future infection burdens}
\[
B_i^0=\sum_{t=31}^{60}I_i^0(t),
\]
\[
B_i^V(a_i)=\sum_{t=31}^{60}I_i^V(t\mid a_i).
\]

\textbf{Expected utility}
\[
EU(a_i)=
\mathbb{E}_{\theta_i\mid C_i}
\left[
B_i^0-B_i^V(a_i)
\right].
\]
\end{minipage}
\hfill
\begin{minipage}[t]{0.47\textwidth}
\textbf{Monte Carlo estimate}
\[
\widehat{EU}(a_i)=
\frac{1}{M}
\sum_{m=1}^{M}
\left[
B_{i,m}^0-B_{i,m}^{V}(a_i)
\right].
\]

\textbf{Decision rule}
\[
\hat{a}
=
\arg\max_{a_i\in A}\widehat{EU}(a_i).
\]

\textbf{Ground-truth optimal action}
\[
a^*
=
\arg\max_{a_i\in A}U^{true}(a_i).
\]
\end{minipage}

\vspace{0.18cm}

\begin{tcolorbox}
\[
\text{Top-1 Accuracy}
=
\mathbb{I}(\hat{a}=a^*),
\qquad
\text{Normalized Utility}
=
\frac{U^{true}(\hat{a})}{U^{true}(a^*)+\varepsilon}.
\]
\end{tcolorbox}

\vspace{0.10cm}

\begin{minipage}[t]{0.47\textwidth}
\[
\textbf{Prediction:}
\quad
\arg\max_i \text{future risk}_i
\]
\end{minipage}
\hfill
\begin{minipage}[t]{0.47\textwidth}
\[
\textbf{Decision:}
\quad
\arg\max_{a_i\in A}
\mathbb{E}\!\left[
B_i^0-B_i^V(a_i)\mid C_i
\right].
\]
\end{minipage}

\vspace{0.10cm}

\[
\text{Action}=\text{allocate vaccine},
\qquad
\text{Utility}=\text{future infection burden reduced}.
\]

\end{document}